% ===================================================================
% Preambel fuer die Druckfassung des FMM-Skripts.
% Wird von scripts/pdf/mdx-to-latex.mjs in den Build-Ordner kopiert.
% Engine: xelatex (defs-fmm.tex braucht \symbf aus unicode-math).
% ===================================================================

\usepackage[a4paper,inner=28mm,outer=30mm,top=25mm,bottom=30mm]{geometry}

% amsmath/mathtools MUESSEN vor unicode-math geladen werden; amssymb darf
% gar nicht mehr dazu (unicode-math bringt die Symbole selbst mit und
% bricht sonst mit "Command \eth already defined" ab).
\usepackage{amsmath}
\usepackage{mathtools}
\usepackage{amsthm}     % defs-fmm.tex: \tagqed braucht \qedhere
\usepackage{centernot}  % defs-fmm.tex: \notimplies braucht \centernot

\usepackage{fontspec}
\usepackage{unicode-math}
\setmainfont{Libertinus Serif}
\setsansfont{Libertinus Sans}
\setmonofont{Libertinus Mono}[Scale=MatchLowercase]
\setmathfont{Libertinus Math}

% shorthands=off ist NICHT optional: babel-german macht " zu einem aktiven
% Zeichen, und die Quellen schliessen ihre „…"-Zitate mit einem geraden
% ASCII-Anfuehrungszeichen. Aus „minimale Krümmung" ist … wird dann
% „minimale Krümmungïst …" ("i -> ï), aus glatt" zu wird glattßu ("z -> ß):
% ein stiller Textschaden mitten im Satz.
\usepackage[ngerman,shorthands=off]{babel}
\usepackage{microtype}
\usepackage{xcolor}
\usepackage{array}
\usepackage{tabularx}
% Gewichtete X-Spalte: \hsize-Faktoren muessen in Summe die Anzahl der
% X-Spalten ergeben. Ohne das bekommt eine Spalte mit lauter "3.1" genauso
% viel Platz wie eine mit drei Zeilen Fliesstext.
% Das \hspace{0pt} ist kein Zierrat: im Flattersatz trennt TeX das ERSTE Wort
% eines Absatzes sonst nicht, und genau die langen deutschen Komposita
% ("Operatornormen") stehen in diesen Tabellen am Zellenanfang.
\newcolumntype{W}[1]{>{\hsize=#1\hsize\raggedright\arraybackslash\hspace{0pt}}X}
\usepackage{booktabs}
\usepackage{enumitem}
\usepackage{needspace}
\usepackage[most]{tcolorbox}
\usepackage{titlesec}
\usepackage{fancyhdr}
\usepackage{listings}
\usepackage[hidelinks,unicode]{hyperref}

% -------------------------------------------------------------------
% Kursmakros. defs-fmm.tex laedt selbst dsfont und benutzt \symbf,
% muss also NACH unicode-math kommen.
% -------------------------------------------------------------------
\input{defs-fmm.tex}

% \left\| … \right\| ist im Skript allgegenwaertig; unicode-math setzt
% \| sonst als Doppelstrich ohne Streckung.
\providecommand{\eps}{\varepsilon}

% Libertinus Math deckt die AMS-Quadrate nicht ab, und amssymb darf neben
% unicode-math nicht geladen werden. Aus Linien gebaut haengen sie an nichts.
\providecommand{\blacksquare}{\mathord{\rule[0.02ex]{0.62ex}{0.62ex}}}
\providecommand{\square}{\mathord{\setlength{\fboxsep}{0pt}\setlength{\fboxrule}{0.4pt}%
  \fbox{\rule{0pt}{0.62ex}\rule{0.62ex}{0pt}}}}

% -------------------------------------------------------------------
% Zu breite abgesetzte Formeln
% -------------------------------------------------------------------
% Im Web scrollt eine ueberlange Formel horizontal; auf Papier ragt sie
% in den Rand — in einem Kasten sogar sichtbar ueber dessen Kante. Diese
% Huelle misst die Formel und verkleinert sie NUR, wenn sie nicht passt.
% Jede Verkleinerung landet als FMM-WIDE-MATH im Log, damit der Bericht
% am Ende sagen kann, wie oft und wie stark eingegriffen wurde.
\newsavebox{\fmmmathbox}
\newcommand{\fmmfit}[1]{%
  \sbox\fmmmathbox{$\displaystyle #1$}%
  \ifdim\wd\fmmmathbox>\linewidth
    \typeout{FMM-WIDE-MATH: \the\wd\fmmmathbox\space > \the\linewidth}%
    \resizebox{\linewidth}{!}{\usebox\fmmmathbox}%
  \else
    \usebox\fmmmathbox
  \fi
}

% -------------------------------------------------------------------
% Farben
% -------------------------------------------------------------------
\definecolor{fmmblue}{HTML}{2563EB}
\definecolor{fmmgreen}{HTML}{15803D}
\definecolor{fmmamber}{HTML}{B45309}
\definecolor{fmmviolet}{HTML}{6D28D9}
\definecolor{fmmslate}{HTML}{475569}
\definecolor{fmmgray}{HTML}{64748B}

% Die Farbmakros aus defs-fmm.tex zeigen im Druck auf die gedeckten
% Varianten — reines green/red/blue ist auf Papier kaum lesbar.
\definecolor{fmmMarkGreen}{HTML}{15803D}
\definecolor{fmmMarkRed}{HTML}{B91C1C}
\definecolor{fmmMarkBlue}{HTML}{1D4ED8}
\definecolor{fmmMarkOrange}{HTML}{C2410C}
\definecolor{fmmMarkPurple}{HTML}{7E22CE}
\renewcommand{\cgreen}[1]{{\textcolor{fmmMarkGreen}{#1}}}
\renewcommand{\cred}[1]{{\textcolor{fmmMarkRed}{#1}}}
\renewcommand{\cblue}[1]{{\textcolor{fmmMarkBlue}{#1}}}
\renewcommand{\corange}[1]{{\textcolor{fmmMarkOrange}{#1}}}
\renewcommand{\cpurp}[1]{{\textcolor{fmmMarkPurple}{#1}}}

% Die \cb*-Varianten setzen in defs-fmm.tex \bfseries — das ist TEXTmodus
% und bricht in einer Formel ab ("\bfseries invalid in math mode"). Im Web
% faellt das nicht auf, weil MathJax den Schalter still ignoriert. Hier
% muss es der Mathe-Fettschnitt sein.
\renewcommand{\cb}[1]{{\symbf{#1}}}
\renewcommand{\cbgreen}[1]{{\textcolor{fmmMarkGreen}{\symbf{#1}}}}
\renewcommand{\cbred}[1]{{\textcolor{fmmMarkRed}{\symbf{#1}}}}
\renewcommand{\cbpurp}[1]{{\textcolor{fmmMarkPurple}{\symbf{#1}}}}
\renewcommand{\cborange}[1]{{\textcolor{fmmMarkOrange}{\symbf{#1}}}}

% -------------------------------------------------------------------
% Ueberschriften. Die Nummern stehen im Quelltext und werden woertlich
% uebernommen (\section* + \addcontentsline), damit Querverweise im
% Fliesstext ("Satz 3.1.4") garantiert stimmen.
% -------------------------------------------------------------------
\titleformat{\chapter}[display]
  {\normalfont\sffamily\bfseries}
  {\color{fmmgray}\large Kapitel \thechapter}
  {6pt}{\Huge}
\titlespacing*{\chapter}{0pt}{10pt}{28pt}

\titleformat{\section}{\normalfont\sffamily\Large\bfseries}{}{0pt}{}
\titlespacing*{\section}{0pt}{20pt}{8pt}
\titleformat{\subsection}{\normalfont\sffamily\large\bfseries\color{fmmslate}}{}{0pt}{}
\titlespacing*{\subsection}{0pt}{14pt}{5pt}
\titleformat{\subsubsection}{\normalfont\sffamily\normalsize\bfseries\color{fmmslate}}{}{0pt}{}
\titlespacing*{\subsubsection}{0pt}{11pt}{4pt}

\setcounter{secnumdepth}{0}
\setcounter{tocdepth}{1}

% Kolumnentitel: "Kapitel 3 · Spur & Matrixnormen" links, Abschnitt rechts.
% Die Abschnitte sind \section* (die Nummern stehen im Quelltext), setzen
% also keine Marke — das erledigt ein \markright je Abschnitt im Dokument.
\renewcommand{\chaptermark}[1]{\markboth{Kapitel \thechapter\ \textperiodcentered\ #1}{}}

\pagestyle{fancy}
\fancyhf{}
\renewcommand{\headrulewidth}{0.3pt}
\fancyhead[LE]{\small\sffamily\color{fmmgray}\leftmark}
\fancyhead[RO]{\small\sffamily\color{fmmgray}\rightmark}
\fancyfoot[C]{\small\sffamily\color{fmmgray}\thepage}
% ACHTUNG: `plain` NICHT per \fancypagestyle umdefinieren. Das darin uebliche
% \fancyhf{} wirkt GLOBAL und loescht die Felder von `fancy` gleich mit —
% sobald die erste Kapitelseite (die \thispagestyle{plain} setzt) gesetzt ist,
% bleibt der Kolumnentitel im ganzen Buch leer. Das eingebaute `plain` aus
% book.cls tut genau das Richtige und fasst fancyhdr nicht an.

% -------------------------------------------------------------------
% Satz-, Definitions-, Beispielumgebungen
% Argumente: {Farbe}{Art}{Label}
% -------------------------------------------------------------------
% \emergencystretch in den Kaesten: das schmalere Satzbild bringt lange
% deutsche Komposita sonst regelmaessig in Zeilen mit klaffenden Wortabstaenden.
\newtcolorbox{fmmenv}[3]{%
  breakable, enhanced, sharp corners,
  boxrule=0pt, leftrule=2.2pt,
  colback=#1!5!white, colframe=#1!75!black,
  left=8pt, right=8pt, top=6pt, bottom=6pt,
  before skip=9pt, after skip=9pt,
  before upper={\emergencystretch=3em
    \noindent\textbf{\sffamily\color{#1!75!black}#2~#3}\par\smallskip}%
}

% Vertiefung: bewusst anders — grau, gestrichelt, kleiner Satz.
\newtcolorbox{fmmvertiefung}[1]{%
  breakable, enhanced, sharp corners,
  boxrule=0.5pt, colback=fmmslate!4!white, colframe=fmmslate!45!white,
  borderline west={2.2pt}{0pt}{fmmslate!55!white, dashed},
  left=8pt, right=8pt, top=6pt, bottom=6pt,
  before skip=9pt, after skip=9pt,
  fontupper=\small,
  before upper={\emergencystretch=3em
    \noindent\textbf{\sffamily\color{fmmslate}Vertiefung:~#1}\par\smallskip}%
}

% Interaktiv: Kasten um ein Widget samt Aufgabe/Auswertung. Im Druck bleibt
% davon die Prosa plus der Widget-Platzhalter; die Marke sagt, dass hier im
% Web etwas zu bedienen ist. Kernstoff, deshalb normale Schriftgroesse.
\newtcolorbox{fmminteraktiv}[1]{%
  breakable, enhanced, sharp corners,
  boxrule=0.5pt, colback=fmmblue!4!white, colframe=fmmblue!45!white,
  borderline west={2.2pt}{0pt}{fmmblue!65!white},
  left=8pt, right=8pt, top=6pt, bottom=6pt,
  before skip=9pt, after skip=9pt,
  before upper={\emergencystretch=3em
    \noindent\textbf{\sffamily\color{fmmblue!75!black}Interaktiv:~#1}\par\smallskip}%
}

% Quellenzeile unter der Abschnittsueberschrift
\newcommand{\fmmquelle}[1]{{\par\small\color{fmmgray}#1\par\addvspace{4pt}}}

% Platzhalter fuer die interaktiven Widgets, die es im Druck nicht gibt.
% Bewusst eine schmale Marke statt eines Kastens: sie soll die Stelle
% markieren, nicht mit den Satz- und Beispielkaesten um Aufmerksamkeit
% konkurrieren. Das Argument (der Komponentenname) bleibt ungenutzt — im
% Druck ist er eine Implementierungsinterna ohne Nutzen fuer die Lesenden.
\newcommand{\fmmwidget}[1]{%
  \par\addvspace{9pt}%
  \noindent\textcolor{fmmgray!40}{\rule{\linewidth}{0.4pt}}\par\vspace{-5pt}%
  {\centering\small\sffamily\color{fmmgray}%
   Interaktives Element — nur in der Web-Fassung\par}%
  \vspace{-3pt}\noindent\textcolor{fmmgray!40}{\rule{\linewidth}{0.4pt}}%
  \par\addvspace{9pt}%
}

% -------------------------------------------------------------------
% Beweise mit Schritten und Begruendungen
% -------------------------------------------------------------------
% Eigenes QED-Zeichen: \square haengt an der Symbolabdeckung der Mathefont,
% ein Rahmen aus Linien haengt an gar nichts.
\newcommand{\fmmqed}{{\setlength{\fboxsep}{0pt}\setlength{\fboxrule}{0.4pt}%
  \fbox{\rule{0pt}{0.62ex}\rule{0.62ex}{0pt}}}}

\newenvironment{fmmproof}[1][qed]{%
  \def\fmmproofqed{#1}%
  \par\addvspace{7pt}%
  \needspace{4\baselineskip}% "Beweis." darf nicht allein am Seitenfuss stehen
  \noindent\textit{Beweis.}\nobreak\par\nobreak\vspace{-3pt}%
  \begin{enumerate}[
      label={\scriptsize\sffamily\color{fmmgray}(\arabic*)},
      leftmargin=*, labelsep=0.7em, align=left,
      topsep=4pt, itemsep=5pt, parsep=2pt]%
}{%
  \end{enumerate}%
  \ifx\fmmproofqed\fmmnoqed\else\nobreak\hfill\fmmqed\par\fi
  \addvspace{7pt}%
}
\def\fmmnoqed{noqed}

% Begruendung eines Beweisschritts
\newcommand{\fmmwhy}[1]{%
  \par\vspace{2pt}%
  \noindent{\footnotesize\color{fmmgray}\textit{#1}}\par%
}

% -------------------------------------------------------------------
% Selbsttest
% -------------------------------------------------------------------
\newenvironment{fmmquiz}{%
  \par\addvspace{6pt}%
  \begin{enumerate}[
      label={\scriptsize\sffamily\color{fmmgray}\arabic*.},
      leftmargin=*, labelsep=0.7em, align=left,
      topsep=3pt, itemsep=7pt, parsep=2pt]%
}{%
  \end{enumerate}\addvspace{6pt}%
}
\newcommand{\fmmantwort}[1]{%
  \par\vspace{2pt}%
  \noindent{\footnotesize\color{fmmslate}\textbf{\sffamily Antwort:}~#1}\par%
}
\newcommand{\fmmloesung}[1]{%
  \par\vspace{2pt}%
  \noindent{\footnotesize\color{fmmslate}\textbf{\sffamily Lösung:}~#1}\par%
}

% Aufklappbares "Lösung anzeigen" aus dem Web
\newcommand{\fmmreveal}[1]{%
  \par\vspace{2pt}\noindent{\small\color{fmmgray}\textbf{\sffamily #1}}\par\vspace{1pt}%
}

% -------------------------------------------------------------------
% Codebloecke
% -------------------------------------------------------------------
% verbatim kann nicht umbrechen: eine lange R-Zeile ragt einfach in den
% Rand. listings bricht um und markiert die Fortsetzung.
\lstset{
  basicstyle=\footnotesize\ttfamily,
  breaklines=true,
  breakatwhitespace=false,
  postbreak={\mbox{\textcolor{fmmgray}{$\hookrightarrow$}\space}},
  columns=fullflexible,
  keepspaces=true,
  extendedchars=true,
  showstringspaces=false,
  frame=none,
  backgroundcolor=\color{fmmslate!6},
  xleftmargin=6pt,
  xrightmargin=2pt,
  framexleftmargin=6pt,
  framexrightmargin=6pt,
  aboveskip=8pt,
  belowskip=8pt,
}

% -------------------------------------------------------------------
% Kleinkram
% -------------------------------------------------------------------
\setlength{\parindent}{0pt}
\setlength{\parskip}{0.55\baselineskip}
\setlist{topsep=3pt, itemsep=2pt, parsep=2pt}
\renewcommand{\arraystretch}{1.15}
\emergencystretch=2em
\hbadness=10000
\vbadness=10000
\clubpenalty=8000
\widowpenalty=8000
\raggedbottom
